% GENERATED FILE. Edit canonical lessons, then run npm run generate:books.
% canonical-source-hash: fnv1a64:3c9ee54c18d5263e
% canonical-lessons: FR-C248-maire, FR-C248-elections, FR-C248-vote, FR-C248-politique, FR-C248-economie

\chapter{Mayor, Elections, Vote, Politics, Economy}
\label{ch:fr-mayor-elections-vote-politics-economy}
\hlchaptermodality{248}

\begin{chapteropening}
\textbf{By the end of this chapter:} \emph{I can say a mayor, elections, a vote, politics and the economy.}

Complete the last lesson of chapter 248: I can say a mayor, elections, a vote, politics and the economy.
\end{chapteropening}

\section[le maire]{le maire — a mayor}

\label{lesson:FR-C248-maire}



\begin{quote}
Before the new one: say the French for a motorbike, then the French for a platform.
\end{quote}

\subsection*{You'll want to know: le maire}
\textbf{le maire} — \textquotedblleft{}a mayor\textquotedblright{}.

\begin{grammarlens}[title={What you've built}]
One more word for this chapter.
\end{grammarlens}

\subsection*{Your turn}
\begin{itemize}
\raggedright
  \item \emph{Say it:} \emph{le maire}
  \item \emph{Say it:} \emph{le maire}, once more
  \item \emph{Say it:} say \emph{le maire}
  \item \emph{Recall:} say the French for a motorbike, then the French for a platform, then say \emph{le maire} again
\end{itemize}

\begin{tcolorbox}[breakable,colback=teal!4,colframe=teal!35!black,title={Before you move on}]
What does le maire mean? (\textquotedblleft{}A mayor\textquotedblright{}.) Say it once more.
\end{tcolorbox}
\section[les élections]{les élections — elections}

\label{lesson:FR-C248-elections}



\begin{quote}
Before the new one: say the French for a platform, then the French for a mayor.
\end{quote}

\subsection*{You'll want to know: les élections}
\textbf{les élections} — \textquotedblleft{}elections\textquotedblright{}.

\begin{grammarlens}[title={What you've built}]
One more word for this chapter.
\end{grammarlens}

\subsection*{Your turn}
\begin{itemize}
\raggedright
  \item \emph{Say it:} \emph{les élections}
  \item \emph{Say it:} \emph{les élections}, once more
  \item \emph{Say it:} say \emph{les élections}
  \item \emph{Recall:} say the French for a platform, then the French for a mayor, then say \emph{les élections} again
\end{itemize}

\begin{tcolorbox}[breakable,colback=teal!4,colframe=teal!35!black,title={Before you move on}]
What does les élections mean? (\textquotedblleft{}Elections\textquotedblright{}.) Say it once more.
\end{tcolorbox}
\section[le vote]{le vote — a vote}

\label{lesson:FR-C248-vote}



\begin{quote}
Before the new one: say the French for a mayor, then the French for elections.
\end{quote}

\subsection*{You'll want to know: le vote}
\textbf{le vote} — \textquotedblleft{}a vote\textquotedblright{}.

\begin{grammarlens}[title={What you've built}]
One more word for this chapter.
\end{grammarlens}

\subsection*{Your turn}
\begin{itemize}
\raggedright
  \item \emph{Say it:} \emph{le vote}
  \item \emph{Say it:} \emph{le vote}, once more
  \item \emph{Say it:} say \emph{le vote}
  \item \emph{Recall:} say the French for a mayor, then the French for elections, then say \emph{le vote} again
\end{itemize}

\begin{tcolorbox}[breakable,colback=teal!4,colframe=teal!35!black,title={Before you move on}]
What does le vote mean? (\textquotedblleft{}A vote\textquotedblright{}.) Say it once more.
\end{tcolorbox}
\section[la politique]{la politique — politics}

\label{lesson:FR-C248-politique}



\begin{quote}
Before the new one: say the French for elections, then the French for a vote.
\end{quote}

\subsection*{You'll want to know: la politique}
\textbf{la politique} — \textquotedblleft{}politics\textquotedblright{}.

\begin{grammarlens}[title={What you've built}]
One more word for this chapter.
\end{grammarlens}

\subsection*{Your turn}
\begin{itemize}
\raggedright
  \item \emph{Say it:} \emph{la politique}
  \item \emph{Say it:} \emph{la politique}, once more
  \item \emph{Say it:} say \emph{la politique}
  \item \emph{Recall:} say the French for elections, then the French for a vote, then say \emph{la politique} again
\end{itemize}

\begin{tcolorbox}[breakable,colback=teal!4,colframe=teal!35!black,title={Before you move on}]
What does la politique mean? (\textquotedblleft{}Politics\textquotedblright{}.) Say it once more.
\end{tcolorbox}
\section[l'économie]{l'économie — the economy}

\label{lesson:FR-C248-economie}



\begin{quote}
Before the new one: say the French for a vote, then the French for politics.
\end{quote}

\subsection*{You'll want to know: l'économie}
\textbf{l'économie} — \textquotedblleft{}the economy\textquotedblright{}.

\begin{grammarlens}[title={What you've built}]
One more word for this chapter.
\end{grammarlens}

\subsection*{Your turn}
\begin{itemize}
\raggedright
  \item \emph{Say it:} \emph{l'économie}
  \item \emph{Say it:} \emph{l'économie}, once more
  \item \emph{Say it:} say \emph{l'économie}
  \item \emph{Recall:} say the French for a vote, then the French for politics, then say \emph{l'économie} again
\end{itemize}

\begin{tcolorbox}[breakable,colback=teal!4,colframe=teal!35!black,title={Before you move on}]
What does l'économie mean? (\textquotedblleft{}The economy\textquotedblright{}.) Say it once more.
\end{tcolorbox}
